%------------------------------------------------------------------------------%
\begin{question}
  Represente o vetor
  \(
    \vec{v} = ( 2, 3 )
  \)
  partindo de
  \(
    A = ( 2, -1 )
  \)
\end{question}

%------------------------------------------------------------------------------%
\begin{resolution}{Solução}
\begin{columns}
\column{0.6\linewidth}
  O vetor \(\vec{v} = ( 2, 3 )\) define

  \pause
  \qquad deslocamento de 2 unidades para a direita

  \pause
  \qquad deslocamento de 3 unidades para cima

  \pause
  \bigskip
  Partindo de \(A = (2, -1)\)

  \pause
  \bigskip
  Avançamos para \(( 4, -1)\)

  \pause
    \bigskip
  Depois para \(( 4, 2 )\)
\column{0.4\linewidth}
  \pause
  \includegraphics{E-01-Vetores-ex-01}
\end{columns}
\end{resolution}

%------------------------------------------------------------------------------%
\endinput
